\chapter*{Introdução}
\addcontentsline{toc}{chapter}{Introdução}
\markboth{Introdução}{}

O biodiesel é obtido industrialmente pela transesterificação de óleos vegetais
ou gorduras animais com um álcool de cadeia curta, quase sempre o metanol, na
presença de um catalisador básico dissolvido no meio reacional. A rota é madura
e opera com rendimentos elevados em tempos de reação curtos, mas carrega
dificuldades que são bem conhecidas de quem a pratica. O catalisador homogêneo
não pode ser recuperado ao fim da reação e precisa ser neutralizado, o que
consome ácido e gera sais. A separação da glicerina exige lavagens sucessivas, e
a água empregada nessas lavagens retorna ao processo como efluente. Quando a
matéria-prima contém ácidos graxos livres em teor apreciável, parte do
catalisador é consumida na formação de sabões, que estabilizam emulsões e
dificultam ainda mais a separação das fases.

A substituição do catalisador dissolvido por um sólido resolve a maior parte
desses problemas de uma só vez. O catalisador deixa de contaminar os produtos,
pode ser retido no reator e reutilizado, e a etapa de neutralização desaparece.
Entre os sólidos básicos estudados para essa finalidade, o óxido de cálcio tem
recebido atenção particular por reunir basicidade adequada, baixo custo e
disponibilidade a partir de fontes residuais, como cascas de ovo e conchas
calcinadas. O preço dessa substituição, entretanto, é a introdução de uma
resistência que a rota homogênea não possui: os reagentes precisam alcançar a
superfície ativa por difusão, e só então reagir.

Essa resistência não é desprezível no caso em questão. Óleos vegetais são
líquidos viscosos cujas moléculas de triglicerídeo têm massa molar próxima de
\SI{870}{\gram\per\mole}, e a difusividade mássica correspondente é da ordem de
\SI{1e-11}{\meter\squared\per\second}. Comparada à difusividade de momento do
mesmo meio, essa grandeza é cerca de cem mil vezes menor, o que coloca o sistema
em uma faixa de número de Schmidt incomum na prática de engenharia de reatores.
Em termos operacionais, isso significa que o transporte de reagente até a parede
catalítica pode se tornar a etapa mais lenta do conjunto, independentemente de
quão ativo seja o catalisador empregado.

É nesse ponto que a configuração do reator deixa de ser um detalhe de projeto.
Reatores monolíticos de parede catalítica, formados por um feixe de canais
paralelos de pequeno diâmetro cujas paredes são recobertas por uma camada fina
de catalisador, oferecem razão entre área superficial e volume muito superior à
de um leito fixo convencional, com perda de carga substancialmente menor. A
área adicional, contudo, só se traduz em desempenho quando existe de fato um
gradiente de concentração a ser vencido. Se o sistema estiver sob controle
cinético estrito, a reação é lenta o bastante para que o reagente alcance a
parede sem dificuldade, e a geometria do monólito não cumpre função alguma que
justifique sua complexidade de fabricação.

Decidir em qual desses regimes o sistema opera não é algo que se resolva por
inspeção. A resposta depende simultaneamente da atividade do catalisador, da
espessura e da porosidade da camada depositada, da geometria do canal e da
velocidade de escoamento, e as grandezas envolvidas percorrem várias ordens de
grandeza. A fluidodinâmica computacional permite tratar esse acoplamento de
forma explícita, resolvendo o campo de concentração no interior do canal com a
reação heterogênea imposta como condição de contorno na parede, e extraindo daí
os grupos adimensionais que caracterizam o regime.

No grupo de pesquisa em que este trabalho se insere, a transesterificação
heterogênea com óxido de cálcio vem sendo investigada experimentalmente há
alguns anos. \citeonline{Melo2023} determinou parâmetros cinéticos para a
metanólise de óleo de soja sobre CaO suportado em biochar a \SI{65}{\celsius},
e \citeonline{Fonseca2026} estendeu o estudo à rota etílica com CaO e
CaO/Al$_2$O$_3$, reportando difusividades efetivas no interior de catalisadores
dessa natureza. Esses dois conjuntos de dados fornecem os parâmetros necessários
para alimentar um modelo de reator, e é sobre eles que o presente trabalho se
apoia.

O objetivo geral desta dissertação é avaliar, por meio de simulação
fluidodinâmica computacional, o desempenho de um reator monolítico de parede
catalítica aplicado à transesterificação heterogênea de óleo de soja com óxido
de cálcio, identificando o regime de operação do sistema e as variáveis de
projeto que o governam. São objetivos específicos:

\begin{itemize}
\itemsep2pt
\item construir um modelo de canal representativo do monólito que resolva
      simultaneamente o escoamento e o transporte multicomponente de espécies,
      com a reação heterogênea tratada como condição de contorno de parede;
\item incorporar a resistência difusional interna ao \textit{washcoat} por meio
      de um fator de efetividade, evitando a discretização da camada porosa;
\item verificar o modelo quanto à conservação de massa, à independência de
      malha e ao fechamento do balanço global por espécie;
\item determinar o número de Damköhler de parede e classificar o regime de
      operação do sistema com os parâmetros cinéticos do catalisador
      efetivamente empregado no laboratório;
\item avaliar a influência da espessura do \textit{washcoat}, da temperatura e
      da geometria do canal sobre a conversão por passe;
\item calibrar o modelo contra os dados experimentais obtidos no laboratório.
\end{itemize}

Nesta redação para o Exame de Qualificação, o Capítulo~\ref{cap:revisao}
apresenta uma revisão da literatura sobre a transesterificação heterogênea, os
catalisadores sólidos básicos empregados na reação e a modelagem de reatores de
parede catalítica, com atenção particular aos grupos adimensionais que
caracterizam a competição entre reação e transporte. O Capítulo~\ref{cap:modelagem}
descreve o modelo matemático construído, as hipóteses adotadas, a implementação
numérica e os testes de verificação aplicados. No Capítulo~\ref{cap:resultados}
reúnem-se os resultados preliminares já obtidos, que incluem a validação
hidrodinâmica do modelo, a determinação do regime de operação e a caracterização
da camada-limite de concentração. O Capítulo~\ref{cap:cronograma} apresenta os
resultados esperados para as etapas restantes e o cronograma previsto para sua
execução.
